% 3-4-features.tex
%  Budget: 7pp -- the largest section of chapter 3.
%  Source map: feature convention follows lago_forecasting_2021 / epftoolbox
%  LEAR exactly (configs/features.yaml header states the cross-check).
%  ALL numbers computed 2026-09-05 from build_features() on the benchmark:
%    X shape (2177, 247); price lags 96; exog lags 96; exog D0 48; weekday 7.
%    price_lag_days [1,2,3,7]; exog_lag_days [1,7]; exog_current_day true.
%  Rules: \lr{} every numeral; \lr{—} for em-dash.

\section{مهندسی ویژگی}

این بخش تعریف می‌کند که مدل در لحظه مرجع پیش‌بینی دقیقاً چه چیزی می‌بیند.
محوری‌ترین قید کل پژوهش نیز در همین‌جا اعمال می‌شود.

\subsection*{تنظیم مسئله}

در مبدأ پیش‌بینی، یعنی روز \lr{O}، هدف پیش‌بینی هر \lr{24} قیمت ساعتی روز
\lr{O+1} است. بنابراین هر سطر ماتریس ویژگی متناظر با \emph{یک روز هدف}
است و نه یک ساعت؛ و بردار پاسخ \lr{24} مؤلفه دارد. این ساختار، مستقیماً
از سازوکار حراج روز-پیش می‌آید که هر \lr{24} ساعت را در یک حراج واحد و
پیش از آغاز روز تعیین می‌کند.

ماتریس ویژگی نهایی \lr{2177} سطر و \lr{247} ستون دارد.

\subsection*{قید علّی: خط قرمز پژوهش}

هیچ ویژگی‌ای مجاز نیست از اطلاعاتی استفاده کند که پس از لحظه مرجع منتشر
می‌شود. این قید ساده به نظر می‌رسد و در عمل به‌آسانی نقض می‌شود، زیرا
بسیاری از متغیرها هم به شکل پیش‌بینی و هم به شکل تحقق‌یافته وجود دارند و
نام ستون‌ها این تمایز را همیشه آشکار نمی‌کند.

در این پژوهش، قید مذکور صرفاً یک قرارداد نگارشی نیست؛ با آزمون خودکار
اعمال می‌شود. آزمونی وجود دارد که مقادیر پس از مبدأ را آشفته می‌کند و
انتظار دارد پیش‌بینی‌ها بدون تغییر بمانند؛ هر وابستگی به آینده، این
آزمون را می‌شکند. افزون بر آن، آزمون دیگری دسترسی شبکه را در حین ساخت
ویژگی مسدود می‌کند، تا تفکیک میان «داده ذخیره‌شده» و «واسط زنده» یک ادعا
نماند.

\subsection*{چهار خانواده ویژگی}

\begin{center}
\begin{latin}
\begin{tabular}{lrl}
\hline
family & columns & definition \\
\hline
price lags        & 96 & 24h price vectors of days $-1,-2,-3,-7$ \\
exogenous lags    & 96 & 24h exog\_1, exog\_2 of days $-1,-7$ \\
exogenous day-ahead & 48 & exog\_1, exog\_2 for the target day itself \\
weekday dummies   &  7 & one-hot day of week of the target day \\
\hline
total             & 247 & \\
\hline
\end{tabular}
\end{latin}
\end{center}

\subsubsection*{وقفه‌های قیمت (\lr{96} ستون)}

بردار کامل \lr{24}‌ساعته قیمت برای روزهای \lr{1}، \lr{2}، \lr{3} و
\lr{7} پیش از روز هدف. انتخاب این چهار وقفه دلخواه نیست: وقفه‌های \lr{1}
تا \lr{3} ساختار کوتاه‌مدت را می‌گیرند و وقفه \lr{7} دوره تناوب هفتگی را،
که هر دو در تحلیل خودهمبستگی بخش \lr{3-3-3} دیده شدند. این انتخاب
هم‌زمان با قرارداد ویژگی مرجع \lr{LEAR} در \cite{lago_forecasting_2021}
یکسان است، و همین یکسانی است که مقایسه بخش \lr{4-5-2} را ممکن می‌کند.

\subsubsection*{وقفه‌های متغیرهای برون‌زا (\lr{96} ستون)}

بردارهای \lr{24}‌ساعته \lr{exog\_1} و \lr{exog\_2} برای روزهای \lr{1} و
\lr{7} پیش از روز هدف.

\subsubsection*{متغیرهای برون‌زای خودِ روز هدف (\lr{48} ستون)}

این خانواده به توضیح نیاز دارد، زیرا در نگاه نخست نقض قید علّی به نظر
می‌رسد: چگونه می‌توان از داده روز هدف استفاده کرد، وقتی آن روز هنوز
نیامده است؟

پاسخ در ماهیت این دو متغیر است. \lr{exog\_1} و \lr{exog\_2} مقادیر
تحقق‌یافته نیستند؛ آن‌ها \emph{پیش‌بینی‌های روز-پیش} بار و تولید
تجدیدپذیرند که توسط بهره‌بردار شبکه و پیش از بسته شدن حراج منتشر
می‌شوند. بنابراین در لحظه مرجع در دسترس‌اند و استفاده از آن‌ها مجاز است.

این نکته را باید صریح نوشت، زیرا خواننده‌ای که ستون‌هایی با برچسب روز
هدف می‌بیند، حق دارد به نشت اطلاعات مشکوک شود. تمایز میان «پیش‌بینی
منتشرشده پیش از حراج» و «مقدار تحقق‌یافته پس از روز» تمام تفاوت را
می‌سازد. اگر مقادیر تحقق‌یافته به کار می‌رفتند، خطای گزارش‌شده به‌طور
مصنوعی کاهش می‌یافت و در بهره‌برداری واقعی دست‌نیافتنی می‌بود.

همین قرارداد در پیاده‌سازی مرجع نیز به کار رفته است؛ بنابراین انتخاب
حاضر، هم از نظر علّی درست است و هم هم‌ارزی پروتکل را حفظ می‌کند.

\subsubsection*{متغیرهای مجازی روز هفته (\lr{7} ستون)}

کدگذاری یک-داغ روز هفته روز هدف. این متغیرها الگوی هفتگی مشاهده‌شده در
بخش \lr{3-3-3} را به‌صورت صریح در اختیار مدل‌هایی می‌گذارند که خود
قادر به استخراج آن نیستند.

\subsection*{آنچه عمداً در مجموعه ویژگی نیست}

سه گروه متغیر که در ادبیات به کار می‌روند، در این پژوهش غایب‌اند و دلیل
غیبت هر یک متفاوت است.

\textbf{قیمت سوخت و شاخص‌های مشتق از آن.} قیمت گاز و زغال‌سنگ و
شاخص‌های مبتنی بر آن‌ها زیر پروانه تجاری منتشر می‌شوند و در دسترس این
پژوهش نبودند. این یک قید است و نه یک انتخاب روش‌شناختی؛ در بخش \lr{5-2}
به‌عنوان محدودیت گزارش می‌شود.

\textbf{متغیرهای هواشناسی.} افزودن آن‌ها هم‌ارزی پروتکل با مجموعه‌داده
مرجع را از میان می‌برد، زیرا مرجع دقیقاً سه سری دارد. حفظ امکان مقایسه
بر افزودن متغیر ترجیح داده شد.

\textbf{متغیرهای مبتنی بر خروج واحدها و ظرفیت در دسترس.} این متغیرها
نیازمند دسترسی ثبت‌نام‌شده‌اند و در زمان تولید نتایج در اختیار نبودند.
پیامد این غیبت در بخش \lr{5-4} به‌عنوان مسیر پژوهش آینده مطرح می‌شود.

\subsection*{هدف روزانه}

افزون بر بردار ساعتی، هدف دوم پژوهش میانگین بار پایه روزانه است. این هدف
از دو مسیر مستقل محاسبه می‌شود: مدل‌سازی مستقیم روی میانگین روزانه، و
میانگین‌گیری از \lr{24} پیش‌بینی ساعتی. هر دو مسیر از همان ماتریس ویژگی
تغذیه می‌شوند، تا تفاوت مشاهده‌شده در بخش \lr{4-4} قابل انتساب به مسیر
باشد و نه به تفاوت اطلاعات ورودی.
